% !TeX root = ../../Book3.tex
\chapter{覆盖空间与全纯覆盖}
\label{b3:ch:11}
\BookChapterNumber{b3:ch:11}
\begin{chapterabstract}
局部逆只处理一个原像点附近的映射；覆盖要求同一个目标邻域被全部原像层完整地覆盖。本章从这一条件证明路径与同伦提升，建立映射提升的基本群判别，并说明覆盖如何携带由底空间拉回的复结构。随后以路径类构造万有覆盖，证明甲板群的自由作用、局部离散性及商的性质。指数、幂映射和圆环给出可直接计算的模型，上半平面的自同构则把自由作用与含固定点的作用区分开来。统一化理论还将进一步确定万有覆盖可能具有的共形形状。
\end{chapterabstract}

\begin{learninggoals}
\begin{itemize}
\item\label{b3:item:11-goal-1}准确区分局部同胚、覆盖和分歧映射，能检查均匀覆盖邻域是否包含全部原像层，并理解适当性如何提供有限层控制。
\item\label{b3:item:11-goal-2}掌握路径、同伦与映射提升的构造，利用闭合提升和基本群像判断一个局部选择能否成为全局映射。
\item\label{b3:item:11-goal-3}理解路径类万有覆盖的拓扑与复结构，区分纤维置换和甲板变换，并准确使用函数下降的条件。
\item\label{b3:item:11-goal-4}计算显式覆盖的层数、甲板群与商坐标，识别上半平面自同构的类型，说明自由性和适当不连续性在商构造中的不同作用。
\end{itemize}
\end{learninggoals}

对数的全部分支可以同时放在一个坐标平面中：以 \(w\) 为变量时，它就是坐标函数；以 \(z=e^w\) 为底变量时，相差 \(2\pi\mathrm i\) 整数倍的点落在同一纤维上。沿底空间绕0一周，提升路径从 \(w\) 走到 \(w+2\pi\mathrm i\)。多值性因此可以改写为一张空间上的单值函数及其投影，而路径造成的分支变化则表现为不同层之间的移动。

\ProseParagraphBreak
但并非每个处处存在局部逆的映射都有这种完整结构。如果删去某个原像点，一条提升路径可能走到缺口而无法继续；底空间的小邻域仍有一些局部逆，却不能再被每一层完整覆盖。覆盖定义中的“全部原像”和“每层映满同一邻域”正是为了排除这一现象。路径提升定理的有限分割只负责把已经完整的局部层接起来，不能补上被删掉的层。

\ProseParagraphBreak
\secref{b3:sec:1101}先完成拓扑部分：均匀覆盖邻域给路径提升，同伦提升控制末点，映射提升判别再把所有闭路的要求写成基本群的包含关系。\secref{b3:sec:1102}将这些结论放进复坐标，讨论全纯函数的提升、下降以及适当无分歧映射。\secref{b3:sec:1103}从底曲面的路径类构造万有覆盖，并核对它确实是 Hausdorff、第二可数且单连通的曲面。\secref{b3:sec:1104}用上半平面矩阵和显式商说明，群作用还必须满足哪些条件才能给普通覆盖。

\ProseParagraphBreak
\chapref{b3:ch:09}中，闭路通过解析延拓置换函数芽；这里将把这种换层现象放到覆盖空间上观察。闭路在纤维上产生的置换、覆盖空间自身的甲板变换、用来构造商的全纯自同构群，是需要区分的三种对象。它们在指数和幂映射中配合得很整齐，但一般覆盖可能没有足够多的甲板变换来交换全部纤维点。万有覆盖给出统一研究这些现象的空间，下一章的统一化则进一步确定它的共形形状。

% !TeX root = ../../Book3.tex
\section{覆盖映射}
\label{b3:sec:1101}

多值函数可以放到一张覆盖空间上变成单值。覆盖的关键条件不是仅有局部逆，而是目标的每个小邻域都被完整、互不相交的各层覆盖。这个条件保证路径与同伦可以一直提升。

\subsection{覆盖映射}
\label{b3:subsec:110101}

\begin{definition}\label{b3:def:11-covering}
连通空间之间的连续满射 \(p:X\rightarrow Y\) 称为\term[fugai-yingshe]{覆盖映射}[covering map]，若每个 \(y\in Y\) 有开邻域 \(V\)，使 \(p^{-1}(V)\) 是互不相交开集 \(V_\alpha\) 的并，且每个 \(p|_{V_\alpha}:V_\alpha\rightarrow V\) 为同胚。这些 \(V_\alpha\) 称为层，\(V\) 称为均匀覆盖邻域。
\end{definition}
本章空间均为 Hausdorff、第二可数的曲面或区域，因此局部路径连通，路径提升所需的局部条件成立。纤维 \(p^{-1}(y)\) 离散，其点数在连通底空间上固定，允许无穷。这个点数称为层数；有限层数不意味着覆盖空间紧，底空间还须紧才有这一推论。

\ProseParagraphBreak
一个局部同胚未必是覆盖，因为某层可能只覆盖目标小邻域的一部分。例如 \(e^z:\mathbb C\setminus\{0\}\rightarrow\mathbb C^*\) 满射且导数从不为零，但1附近有一个被删掉中心的局部对数层，无法完整覆盖该邻域。因此它不是覆盖。完整指数 \(e^z:\mathbb C\rightarrow\mathbb C^*\) 才是覆盖。

\ProseParagraphBreak
局部同胚与覆盖的差别可以用层是否完整来判断。每一点有局部逆只说明一层在该点存在；覆盖要求同一个目标小邻域同时容纳所有原像层，并且没有层在邻域内突然结束或逃逸。开区间 \((0,3/2)\) 上的 \(t\mapsto e^{2\pi\mathrm{i}t}\) 是到圆周的满射局部同胚，但不是覆盖：在某些目标点有两层，在某些点只有一层。有限端点截断了一个本来应完整的层。这个实一维例子足以显示定义中的全部原像条件不是冗余的。

\ProseParagraphBreak
覆盖的层数在连通目标上恒定，允许无限。沿任一连接两点的路径提升，可把起始纤维的每一点送到终点纤维；逆路径给逆映射，因而纤维基数相同。这个论证不要求有限层，也不要求目标单连通。不同路径可能给不同纤维双射，闭路造成的置换正是覆盖的 monodromy 资料。

\subsection{路径提升}
\label{b3:subsec:110102}

\begin{theorem}[路径提升]\label{b3:thm:11-path-lifting}
给定覆盖 \(p:X\rightarrow Y\)、连续路径 \(\gamma\) 及 \(x_0\in p^{-1}(\gamma(0))\)，存在唯一连续路径 \(\widetilde\gamma\) 满足 \(p\circ\widetilde\gamma=\gamma\)、\(\widetilde\gamma(0)=x_0\)。
\end{theorem}
\begin{proof}
用均匀覆盖邻域覆盖紧路径像，按 Lebesgue 数把参数分成有限小段，每段像落在一个邻域中。首段使用包含 \(x_0\) 的层的逆映射；它的终点确定下一段所用层，依次得到连续提升。两提升若在一个位置相同，则在附近使用同一层逆而相同；相同位置集合开且闭，且包含0，故由区间连通性处处相同。
\end{proof}
路径每绕一个孔，提升终点可移到纤维的另一个点。起点选择是必要数据；若只给底路径，有限层覆盖通常有多条不同提升。闭路的提升也未必闭合，指数覆盖的单位圆提升从0走到 \(2\pi\mathrm{i}\)。

\subsection{同伦提升}
\label{b3:subsec:110103}

\begin{theorem}[同伦提升]\label{b3:thm:11-homotopy-lifting}
若 \(H:[0,1]^2\rightarrow Y\) 为路径同伦，给定 \(H(0,\cdot)\) 的一个提升，便有唯一提升 \(\widetilde H\) 延伸它。端点固定同伦的提升在给定起点后具有相同终点。
\end{theorem}
\begin{proof}
将参数方形分为有限小格，使每格像落在均匀覆盖邻域。已知第一侧提升后，逐格用包含一个已知顶点的层逆定义格上提升。公共边的两个定义是同一底路径的提升，且在端点相同，路径提升唯一性使它们一致，故所有格拼成连续提升。唯一性同样由局部层逆和连通性得到。若端点固定，提升终点沿 \(s\) 形成落在一个离散纤维中的连续路径，因此恒定。
\end{proof}
这使底空间基本群作用在纤维上：沿闭路提升取末点。作用只依赖同伦类，但它未必与甲板群相同；甲板群是整张覆盖的对称，而纤维作用是路径传递。

\begin{theorem}[映射提升的判别]\label{b3:thm:11-map-lifting}
设 \(p:X\to Y\) 是连通曲面之间的覆盖，\(F:Z\to Y\) 连续，\(Z\) 是连通曲面。选 \(z_0\in Z\)、\(x_0\in X\)，使 \(p(x_0)=F(z_0)=y_0\)。存在满足 \(\widetilde F(z_0)=x_0\) 的提升，当且仅当
\[
 F_*\pi_1(Z,z_0)\subset p_*\pi_1(X,x_0).
\]
满足基点条件的提升唯一；若各曲面带复结构，\(p,F\) 全纯，则提升全纯。
\end{theorem}
\begin{proof}
若有提升，闭路先在 \(Z\) 映到 \(X\) 再投到 \(Y\)，故群包含必要。反向对 \(z\) 选从 \(z_0\) 到 \(z\) 的路径 \(\gamma\)，把 \(F\circ\gamma\) 从 \(x_0\) 提升，以终点定义 \(\widetilde F(z)\)。两条路径拼成闭路，其像类属于 \(p_*\pi_1(X)\)。因此它与某条 \(X\) 闭路的投影同伦，利用同伦提升，像闭路的提升也闭合。两条原路径的提升终点相同，定义不依赖选择。每个 \(z\) 附近选小连通邻域，使其像落在覆盖均匀邻域；这时提升等于固定局部逆与 \(F\) 的复合，故连续。两提升相等的点集在局部逆图内开，其不相等的点集也开，连通性给唯一。全纯情形的同一局部复合表明提升全纯。
\end{proof}
特别地，单连通源的群为0，所以总可提升。第十三章把圆盘上的全纯映射提升到目标曲面的圆盘覆盖时，用的正是这个特例。源仅为一个圆环时没有这样的自动结论：映射在基本圈上绕行的次数必须属于目标覆盖对应的子群。

\ProseParagraphBreak
例如 \(p(w)=w^m:\mathbb C^*\to\mathbb C^*\)，其基本群像是 \(m\mathbb Z\)。一个无零全纯函数 \(F:U\to\mathbb C^*\) 提升到这张覆盖，恰是找到 \(h^m=F\)；判别条件就是所有像闭路的绕数被 \(m\) 整除。这与第十章根的周期条件相同，两个观点分别把它写成群包含和对数导数积分。

\subsection{甲板变换}
\label{b3:subsec:110104}

\begin{definition}\label{b3:def:11-deck}
满足 \(p\circ T=p\) 的覆盖空间自同胚称为\term[jiaban-bianhuan]{甲板变换}[deck transformation]，全部甲板变换组成群 \(\operatorname{Deck}(p)\)。
\end{definition}
\begin{proposition}\label{b3:prop:11-deck-free}
连通覆盖的甲板变换若固定一点，便是恒等。每点有小邻域与其任何非恒等甲板像不交。
\end{proposition}
\begin{proof}
连接固定点到任意点的路径，经 \(T\) 变换后仍是同一底路径的提升且起点相同，唯一性使末点不变，所以 \(T\) 恒等。对均匀覆盖邻域中的一层 \(W\)，若 \(T(W)\) 与 \(W\) 相交，它们同为该邻域的层，故相同；层上投影单射使 \(T\) 在 \(W\) 恒等，再由第一结论为全域恒等。
\end{proof}
任意覆盖的甲板群不必在纤维上传递。若它确实传递，称该覆盖为正规覆盖；万有覆盖将满足这一性质。不要把“局部各层相同”直接解释为存在交换任意两层的全局甲板变换。

% !TeX root = ../../Book3.tex
\section{全纯覆盖}
\label{b3:sec:1102}

全纯覆盖把上节的拓扑条件与复坐标结合。它处处局部双全纯；与存在临界点的分歧映射相比，解析函数的提升因此没有局部角度折叠。

\subsection{全纯覆盖}
\label{b3:subsec:110201}

\term[riemann-qumian]{Riemann 曲面}[Riemann surface]是 Hausdorff、第二可数且具有复平面局部坐标的连通曲面，各坐标变换全纯且导数非零。曲面间映射全纯，指其坐标表达全纯。若一个拓扑覆盖投影到 Riemann 曲面，则在每一层拉回底坐标，给覆盖空间唯一的 Riemann 曲面结构，使投影局部双全纯。坐标交叠的变换直接是底空间的坐标变换，所以相容性完整成立。

\ProseParagraphBreak
在这个结构下甲板变换自动全纯：在一层坐标中它把相同底坐标传到另一层，坐标表达为恒等。全纯函数 \(f\) 的提升 \(f\circ p\) 也全纯。反向，一个甲板不变全纯函数可按底点的任意提升定义底函数，各层逆保证它局部全纯；这就是下降。

\subsection{指数映射}
\label{b3:subsec:110202}

\begin{proposition}\label{b3:prop:11-exponential-cover}
指数 \(p(w)=e^w\) 是 \(\mathbb C\rightarrow\mathbb C^*\) 的万有覆盖，其甲板群为 \(w\mapsto w+2\pi\mathrm{i}n\)、\(n\in\mathbb Z\)。
\end{proposition}
\begin{proof}
在任意目标点的小无零圆盘上取全纯对数 \(L\)。全部原像层为 \(L(V)+2\pi\mathrm{i}n\)，互不交，指数在每层有逆 \(L+2\pi\mathrm{i}n\)，所以是覆盖。源平面单连通，故为万有覆盖。甲板变换满足 \(e^{T(w)-w}=1\)，连续性使 \(T(w)-w\) 是固定的 \(2\pi\mathrm{i}n\)。反向这些平移显然是甲板变换。
\end{proof}
对数在覆盖空间上就是坐标 \(w\)。它不下降为底空间单值函数，因为甲板平移改变它；\(e^{mw}\) 在 \(m\in\mathbb Z\) 时不变，下降为 \(z^m\)。这种下降判别把分支周期转成群不变性。

\subsection{去心区域的覆盖}
\label{b3:subsec:110203}

圆环 \(A(r,R)\) 的指数覆盖限制在 \(\log r<\operatorname{Re}w<\log R\)。去心盘的覆盖限制在左半平面。它们的源域都单连通，故是各自万有覆盖。把左半平面用分式变换映成圆盘，可得到圆盘到去心盘的显式覆盖，但这个映射不是单射；其甲板变换在圆盘坐标中变为抛物型自同构。

\ProseParagraphBreak
有限覆盖 \(w\mapsto w^m:\mathbb C^*\rightarrow\mathbb C^*\) 的甲板群为 \(w\mapsto\zeta w\)、\(\zeta^m=1\)。单位正圈提升只转过 \(2\pi/m\)，走 \(m\) 次才闭合。它的源仍不单连通，因此不是万有覆盖；有限层与万有性是不同条件。

\begin{proposition}[圆环的有限层覆盖]\label{b3:prop:11-annulus-power-cover}
\(z\mapsto z^m\) 把 \(A(r,R)\) 作为 \(m\) 层全纯覆盖映到 \(A(r^m,R^m)\)，目标模量是源模量的 \(m\) 倍。
\end{proposition}
\begin{proof}
目标每一点有 \(m\) 个根，模均落在 \((r,R)\)。导数 \(mz^{m-1}\) 在圆环非零；选足够小的目标圆盘可取 \(m\) 个不同根分支，且它们的像互不相交、覆盖全部原像，故满足均匀覆盖定义。模量由对数半径比计算，即 \(\log(R^m/r^m)/(2\pi)=m\log(R/r)/(2\pi)\)。
\end{proof}
有限层覆盖不保环域模量，因为它并非一一对应的共形等价。第八章模量不变性针对双全纯映射，不能把“全纯且导数非零”自动升级成那个假设。通过条带 \(\log r<\operatorname{Re}w<\log R\) 的指数覆盖看，同一事实是甲板平移周期由 \(2\pi\mathrm{i}\) 的格点变为一个有限指数子群。

\ProseParagraphBreak
对于去心盘，指数覆盖使用左半平面，去心平面则使用整个平面。两者都只有一个可见的去心点，却给不同万有覆盖类型：前者的外边界提供额外限制，后者没有外边界。把所有“去心域”只画成一个孔，容易忽略这个全局差别。

\subsection{与分歧覆盖的区别}
\label{b3:subsec:110204}

映射 \(z\mapsto z^m\) 在含0的平面上有局部次数 \(m\)，0处导数为零，不能作为覆盖。删去分歧点及其值后，才成为前一小节的有限覆盖。

\begin{proposition}\label{b3:prop:11-proper-unbranched}
若非恒常全纯映射 \(f:X\rightarrow Y\) 适当且无临界点，即紧集原像紧，则它是有限层全纯覆盖；若 \(Y\) 连通，层数固定。
\end{proposition}
\begin{proof}
无临界点给局部双全纯，故像开。紧集原像紧使像闭：若 \(f(x_n)\rightarrow y\)，这些像点连同极限构成紧集，原像紧使 \(x_n\) 有子列趋于 \(x\)，故 \(f(x)=y\)。像非空，所以满射。每个纤维既离散又紧，因此有限。围绕有限原像点取互不交逆函数邻域；缩小底邻域后不可能还有额外原像，否则一列额外点像趋于 \(y\)，紧集原像紧给极限落入某个既选邻域，矛盾。故底邻域均匀覆盖，有限层数由连通性固定。
\end{proof}
满足紧集原像紧的映射称为\term[shidang-yingshe]{适当映射}[proper map]；这里的条件是一项全局紧性要求。多项式去掉临界值及其全部原像后满足这些条件，给出典型有限覆盖。

\ProseParagraphBreak
多项式 \(P\) 的适当性可直接从最高次项证明。\(|z|\) 足够大时 \(|P(z)|\ge c|z|^d\)，所以紧目标集的原像有界；连续性又使它闭，故紧。删去有限临界值及其全部原像，\(P\) 就成为 \(d\) 层无分歧覆盖。以 \(P(z)=z^3-3z\) 为例，临界值为 \(\pm2\)，删掉它们后每个目标点有三个不同根；小目标盘内全部根分支由局部逆取得，适当性排除第四种逃向无穷的局部层。重新补回临界值后，局部 \(z^2\) 形状在相应纤维出现，完整映射是分歧覆盖。

\ProseParagraphBreak
另一方面，\(e^z\) 不适当：值1的原像 \(2\pi\mathrm{i}\mathbb Z\) 无界。它仍是覆盖，因为指数的周期让全部无穷层有同一个均匀邻域。适当性是取得有限层覆盖的充分工具，不是覆盖的必要条件。应分别检查导数、适当性、满射与完整层分解，不能用其中一项替代其余。

% !TeX root = ../../Book3.tex
\section{万有覆盖}
\label{b3:sec:1103}

万有覆盖用路径类制造一张使所有闭路都能收缩的曲面。与最大解析延拓不同，这个构造只取决于底空间拓扑，随后才拉回复结构。

\subsection{万有覆盖}
\label{b3:subsec:110301}

固定 \(a\in Y\)。令 \(\widetilde Y\) 为从 \(a\) 出发的路径端点固定同伦类，投影取终点。对终点 \(y\) 的一个小坐标圆盘 \(V\)，以及一条到 \(y\) 的路径 \(\gamma\)，把 \([\gamma*\eta]\)、其中 \(\eta\) 在 \(V\) 从 \(y\) 到变动终点，作为一个坐标片。圆盘单连通使类只依赖终点，故这个片与 \(V\) 同胚。两片若有同一点，则在共同更小盘的部分相同，否则互不交。

下面构造\term[wanyou-fugai]{万有覆盖}[universal covering]，其源空间单连通。

\begin{theorem}[万有覆盖]\label{b3:thm:11-universal-cover}
每张连通 Riemann 曲面具有单连通的全纯覆盖 \(p:\widetilde Y\rightarrow Y\)，给定一个纤维点后它在覆盖同构意义下唯一。
\end{theorem}
\begin{proof}
上述片构成拓扑，投影在每片同胚，故是覆盖。路径类可沿其代表逐步提升，说明 \(\widetilde Y\) 路径连通。端点不同的两类由不交底盘分开；端点相同而类不同，则同一小盘的两层不交，所以空间 Hausdorff。底曲面有可数小盘基，有限盘链及起始类确定可达层，因此得到可数基。

给覆盖空间基点闭路 \(\beta\)，其投影为 \(\alpha\)。路径提升定义使 \(\beta\) 的终点类恰为 \([\alpha]\)；闭合意味着 \([\alpha]\) 为平凡类。把底闭路的收缩同伦提升，得到 \(\beta\) 的收缩，故源单连通。拉回坐标给全纯结构。若另有单连通覆盖，沿同一路径从选定点提升取末点；同伦提升使它依类良定义，并在每小片是底坐标的恒等对应。反向同样构造，两复合因路径提升唯一性为恒等，得到唯一同构。
\end{proof}

\subsection{基本群作用}
\label{b3:subsec:110302}

基点闭路类 \([\alpha]\in\pi_1(Y,a)\) 对路径类作 \([\gamma]\mapsto[\alpha*\gamma]\)。同伦相容性与逆路径给这是甲板变换。任意甲板变换由它对基点纤维的像决定，而该像正是一个闭路类，所以这种作用给甲板群与基本群的对应。

\ProseParagraphBreak
拼接顺序采用先走左因子再走右因子，群表示的写法随这个约定固定；如果文献用另一种路径乘法，可能出现反群。几何结论是相同的：基本群记录万有覆盖各层，闭路提升末点记录该类所对应的甲板平移。

\ProseParagraphBreak
连通覆盖的纤维置换作用是传递的。任意两点 \(x_1,x_2\) 在同一纤维，可在源曲面中选连接路径；其投影是目标闭路，从 \(x_1\) 提升后终点为 \(x_2\)。但传递性不等于存在把任意纤维点交换的甲板变换。甲板变换必须同时与整张覆盖相容，不是一个纤维上的任意置换。万有覆盖的甲板群才有正文所证的自由传递性。

\ProseParagraphBreak
例如有限覆盖的基本群子群为 \(H\)，闭路作用可以在余类上描述，甲板群则只由保持 \(H\) 的整体对称产生。如果 \(H\) 不是正规子群，一张连通覆盖可能有很少的甲板变换，尽管所有纤维点仍由闭路作用互相可达。这说明解析延拓产生的根置换，与某张曲面自身的全局对称，也须分别记录；二者在高度对称的 \(z^m\) 例子中恰好紧密对应，不应把该例的特殊性写成一般事实。

\subsection{甲板群}
\label{b3:subsec:110303}

\begin{proposition}\label{b3:prop:11-universal-deck}
万有覆盖的甲板群在每个纤维上传递且自由，底曲面与商 \(\widetilde Y/\operatorname{Deck}(p)\) 自然同构。
\end{proposition}
\begin{proof}
两个同终点路径类 \([\gamma],[\eta]\) 的差可由闭路 \(\eta*\gamma^{-1}\) 补偿，前置这条闭路将第一类变成第二类。故传递；自由性已由甲板固定点命题证明。投影每个纤维恰是一个轨道，诱导商到 \(Y\) 的双射。小层邻域上投影同胚，商映射同样局部同胚，因此双射为同胚；坐标相同使它双全纯。
\end{proof}
反向，若一组全纯自同构自由且适当地不连续作用于曲面，商具有拉回自局部不交邻域的复结构，投影为覆盖。自由性不能删除，否则商坐标将包含分歧；局部不交条件不能删除，否则商可能不再 Hausdorff。

\begin{proposition}[甲板作用的局部离散性]\label{b3:prop:11-deck-proper}
万有覆盖的甲板群在源曲面上自由作用。每一点有小邻域 \(V\)，使 \(T(V)\cap V=\varnothing\) 对所有非恒等甲板变换成立；任意两个紧集 \(K,L\) 之间只有有限个甲板变换满足 \(T(K)\cap L\ne\varnothing\)。
\end{proposition}
\begin{proof}
取均匀覆盖邻域的一层 \(V\)。若 \(T(V)\) 与 \(V\) 相交，两者作为层便相同；由于投影在该层一一对应，\(T\) 在 \(V\) 为恒等，唯一提升给全局恒等。若紧集结论失败，取不同 \(T_n\) 与 \(x_n\in K\)，使 \(T_nx_n\in L\)。子列上 \(x_n\to x\)、\(T_nx_n\to y\)，且 \(p(x)=p(y)\)。取该目标点的一个均匀邻域，令含 \(x,y\) 的层为 \(V,W\)。充分大时 \(x_n\in V\)、\(T_nx_n\in W\)，所以所有 \(T_n\) 都把 \(V\) 送到 \(W\)。投影一一性使它们在 \(V\) 相同，从而全局相同，矛盾。
\end{proof}
这给商的局部坐标：在 \(V\) 中没有不同点被甲板群识别，所以商映射在那里是同胚；目标曲面恰由这些图拼成。紧集有限相交也是后面在双曲圆盘中求最短甲板距离、证明商度量完备的重要事实。不能只说群“离散”而省略到底是在变换空间中离散还是轨道在曲面中局部离散。

\subsection{Riemann 曲面观点}
\label{b3:subsec:110304}

函数芽曲面若其投影确为覆盖，可以由底空间基本群的作用描述；一般最大芽空间的局部双全纯投影未必完整覆盖，因为某些路径遇到奇点无法继续。万有覆盖则对每条底路径都能提升。二者的关系是：若一个初始芽能沿底空间全部路径延拓，拉到万有覆盖后由单值性定理成为单值全纯函数，下降到底空间当且仅当它对全部甲板变换不变。

\ProseParagraphBreak
这给对数与幂函数一个统一位置，也为一般区域的 Poincaré 度量铺路。局部坐标随甲板变换而变，但圆盘的内在密度保持不变，所以可以下降。要知道万有覆盖是否圆盘，还需要统一化输入，不能由拓扑万有覆盖存在性直接推出。

\ProseParagraphBreak
拓扑覆盖携带复结构的过程具有唯一性。每一层用目标坐标与投影复合成源坐标，层间转换正是目标转换，故全纯。若另一复结构也使投影局部双全纯，两结构在每个层上都由同一个目标坐标拉回，必相容。这意味着覆盖问题的主要新资料在全局路径与纤维，局部解析结构已经被目标决定。分歧覆盖则还需补充局部重数，不能通过普通覆盖的坐标拉回直接越过分歧点。

\ProseParagraphBreak
读者可以把本章三个构造连在一起核验：先由路径同伦类形成万有覆盖，再以层坐标给它复结构，最后以闭路拼接给甲板变换。任一构造都不使用统一化；只有在下一章把这个已建曲面识别为球面、平面或圆盘时，才加入深层分析输入。先构造对象再识别模型，使后面度量的定义不依赖一个尚未写清的覆盖空间。

% !TeX root = ../../Book3.tex
\section{上半平面与模群预告}
\label{b3:sec:1104}

圆盘与上半平面是同一个共形模型。上半平面的自同构用实矩阵表示，既方便分类甲板变换，也给模群与模形式提供几何入口。

\subsection{上半平面}
\label{b3:subsec:110401}

上半平面记为 \(\mathbb H=\{z:\operatorname{Im}z>0\}\)。Cayley 变换 \(C(z)=(z-\mathrm{i})/(z+\mathrm{i})\) 将它双全纯映到圆盘，逆为 \(\mathrm{i}(1+w)/(1-w)\)。实轴加无穷远对应单位圆周，内部对称可在两种模型间共轭传递。

\ProseParagraphBreak
无穷远是球面边界点，不是半平面内部点。平移 \(z\mapsto z+1\) 没有内部固定点，却在边界固定无穷远；缩放 \(z\mapsto az\)、\(a>0\)、\(a\ne1\) 固定边界的0与无穷远。两种无固定点运动的商具有不同端部几何。

\subsection{Möbius 变换}
\label{b3:subsec:110402}

实系数分式 \(T(z)=(az+b)/(cz+d)\)、\(ad-bc>0\)，满足
\[
 \operatorname{Im}T(z)=\frac{(ad-bc)\operatorname{Im}z}{|cz+d|^2}>0.
\]
所以它是半平面自同构；反向由圆盘自同构分类与 Cayley 共轭得到所有自同构都如此。将矩阵整体除以行列式平方根，可令行列式为1；矩阵与其负矩阵给相同分式。

\ProseParagraphBreak
按实扩充轴上的固定点数，非恒等自同构分为椭圆型、抛物型、双曲型：分别有内部固定点、一个重边界固定点、两个不同边界固定点。对行列式1矩阵，判别式为 \((a+d)^2-4\)，故迹绝对值分别小于、等于、大于2。甲板变换自由，因而不能是非平凡椭圆型。

\subsection{\texorpdfstring{\(\mathrm{PSL}_2(\mathbb R)\)}{PSL2(R)}}
\label{b3:subsec:110403}

群 \(\operatorname{PSL}_2(\mathbb R)=\operatorname{SL}_2(\mathbb R)/\{\pm I\}\) 正是保定向半平面全纯自同构群。后面的双曲度量元 \(|\mathrm{d}z|/\operatorname{Im}z\) 在它作用下不变，可由上式与 \(T'(z)=1/(cz+d)^2\) 立即验证。

\ProseParagraphBreak
群的离散性是关于变换参数的拓扑，适当不连续性是关于在半平面的作用；对这里的离散子群，作用的局部性质可由矩阵紧性与内部点控制证明。要形成普通无分歧覆盖，还必须无挠，排除椭圆固定点。只写一个 \(\operatorname{PSL}_2(\mathbb R)\) 子群，并不足以自动得到光滑商曲面。

\ProseParagraphBreak
实 Möbius 变换的不动点类型可由二次方程核验。把矩阵归一化为行列式1，不动点方程的判别式为 \((a+d)^2-4\)。当其负时，两根共轭且一根在上半平面，是椭圆型；当其为0且不是恒等变换，只有一个实边界不动点，是抛物型；当其正时，有两个不同实边界不动点，是双曲型。\(\operatorname{tr}^2\) 而非迹本身才与 \(\mathrm{PSL}_2\) 中矩阵正负代表无关。

\ProseParagraphBreak
标准形分别是圆盘内旋转、上半平面平移和正伸缩。前者有内部不动点，后两者没有。覆盖的甲板群须自由，故不能含非平凡椭圆元。\(\mathrm{PSL}_2(\mathbb Z)\) 包含椭圆元，并不满足这个自由条件；完整模群商有特殊稳定子点，不能直接称为通常无分歧曲面覆盖。若转向没有稳定子的子群，才能按前面覆盖理论形成无分歧商；群的离散性和作用的自由性是不同条件。

\subsection{\texorpdfstring{\(\mathrm{PSL}_2(\mathbb Z)\)}{PSL2(Z)}}
\label{b3:subsec:110404}

模群 \(\operatorname{PSL}_2(\mathbb Z)\) 由整数矩阵给出，经典生成元为 \(T(z)=z+1\)、\(S(z)=-1/z\)。欧氏算法消去整数矩阵的下左元素，可将任意行列式1整数矩阵写成这两类操作的乘积，故它们生成模群。

\ProseParagraphBreak
\(S\) 固定 \(\mathrm{i}\)，\(ST\) 固定 \((-1+\sqrt3\,\mathrm{i})/2\)。因此整个模群商具有椭圆点，商投影不是处处无分歧覆盖。通常基本区域 \(|\operatorname{Re}z|\le1/2,|z|\ge1\) 的两竖边以 \(T\) 识别，圆弧边以 \(S\) 识别；无穷远对应尖点。模形式将要求全纯函数在变换下按 \((cz+d)^k\) 的权变换，并另检验尖点展开。这里的几何预告不先假定模形式的维数或傅里叶展开定理。


\begin{chaptersummary}
覆盖以完整的局部层保证路径和同伦始终能够提升；起点一经指定，提升便唯一。映射能否提升取决于其基本群像是否落在覆盖对应的子群中，根的周期条件由此得到几何解释。全纯结构可沿层坐标唯一拉回，函数下降则要求每条纤维上的值相同；只有正规覆盖才能把这个条件改写成甲板不变性。万有覆盖由路径类构成，它的甲板群在纤维上自由传递，底曲面是相应的商。反向从群作用构造覆盖时，自由性排除局部分歧，适当不连续性保证不同轨道能被分开。有限层覆盖的紧性与适当性、指数覆盖的无穷层、椭圆元的内部固定点，分别提醒我们不能省略这些全局条件。本章证明覆盖的存在和基本性质，并计算若干显式共形模型；一般统一化是下一阶段的分析问题。
\end{chaptersummary}

\BookExercises

% 本章八题自拟或由旧题实质改编，未引用未核实的外部原题。
% 各题独立使用正文工具；题8以正文已算出的缩放商圆环模型为起点，新增反演及分歧商训练。
% 训练增量：开路径末点、多个孔的提升限制、非离散周期、有限覆盖的proper性、非正规六层覆盖、函数与微分下降、生成元与群的自由性、商上再取商。
\Needspace{6\baselineskip}
\begin{exercise}\label{b3:ex:11-1}
设 \(\gamma,\eta\) 是 \(\mathbb C^*\) 中从 \(a\) 到 \(b\) 的两条分段光滑路径，记
\[
 n=\operatorname{Ind}(\gamma*\eta^{-1},0).
\]
将两条路径提升到指数覆盖，并指定同一起点 \(u_0\)、\(e^{u_0}=a\)；再将它们提升到 \(p_m(v)=v^m\)、\(m\ge2\)，指定同一起点 \(v_0\)、\(v_0^m=a\)。分别求两种覆盖下末点之差或之比，说明有限幂覆盖保留了哪些绕数信息。应用于
\[
 \gamma(t)=(1+t)e^{4\pi\mathrm i t},\quad
 \eta(t)=1+t,\quad u_0=0,\quad v_0=1,
\]
并指出两条不同伦路径在哪个 \(m\ge2\) 下仍有相同的提升末点。
\end{exercise}
\begin{solution}
指数提升在每段小邻域中由一个对数分支给出，因此可逐段微分并积分，得到
\[
 \widetilde\gamma(1)-u_0=\int_\gamma\frac{\mathrm dz}{z},\quad
 \widetilde\eta(1)-u_0=\int_\eta\frac{\mathrm dz}{z}.
\]
两式相减，即有
\[
 \widetilde\gamma(1)-\widetilde\eta(1)=2\pi\mathrm i n.
\]
所以两条指数提升的末点相同当且仅当 \(n=0\)。

利用已经求得的指数提升，可写幂覆盖中的提升为
\[
 \widehat\gamma(t)=v_0
       \exp\left(\frac{\widetilde\gamma(t)-u_0}{m}\right).
\]
它的 \(m\) 次方是 \(\gamma(t)\)，起点为 \(v_0\)，故由路径提升唯一性就是所需提升；\(\eta\) 同样处理。于是
\[
 \frac{\widehat\gamma(1)}{\widehat\eta(1)}
 =e^{2\pi\mathrm i n/m}.
\]
幂覆盖的末点只记录绕数模 \(m\)，两末点相同等价于 \(m\mid n\)，这比 \(n=0\) 弱。

具体两条路径的指数提升分别为 \(\log(1+t)+4\pi\mathrm i t\) 与 \(\log(1+t)\)，末点分别为 \(\log2+4\pi\mathrm i\) 与 \(\log2\)。幂覆盖末点为
\[
 2^{1/m}e^{4\pi\mathrm i/m},\quad 2^{1/m}.
\]
其相等条件是 \(m\mid2\)，所以在 \(m\ge2\) 中恰为 \(m=2\)。两路径的差闭路绕数为2，故它们在 \(\mathbb C^*\) 中不端点固定同伦；平方根覆盖的末点仍不能分辨这两个类。有限覆盖的闭合提升条件不能当作底路径可缩的判别。
\end{solution}

\Needspace{6\baselineskip}
\begin{exercise}\label{b3:ex:11-2}
设 \(m\ge2\)、\(k,\ell\in\mathbb Z\)，在 \(U=\mathbb C\setminus\{0,1\}\) 上令 \(F(z)=z^k(z-1)^\ell\)。对正整数 \(n\)，取
\[
 U_n=\mathbb C\setminus\bigl(\{0\}\cup\{\zeta:\zeta^n=1\}\bigr),
 \quad q_n(\zeta)=\zeta^n.
\]
求 \(F\circ q_n\) 能提升到覆盖 \(p_m(w)=w^m:\mathbb C^*\to\mathbb C^*\) 的充要条件；若存在这样的 \(n\)，求最小值并写出全部提升。分别讨论 \((m,k,\ell)=(6,4,3)\) 与 \((6,4,6)\)，说明为什么只检查绕0的闭路会给出错误判断。
\end{exercise}
\begin{solution}
有
\[
 (F\circ q_n)(\zeta)=\zeta^{kn}(\zeta^n-1)^\ell.
\]
若存在提升 \(h\)，则 \(h^m=F\circ q_n\)，每条像闭路的绕数必须被 \(m\) 整除。在0附近的小正圆上，右边的绕数为 \(kn\)。在任意 \(n\) 次单位根 \(\xi\) 附近，\(\zeta^n-1\) 的零点是单的，因为 \(n\xi^{n-1}\ne0\)，而 \(\zeta^{kn}\) 无零无极点，所以像绕数为 \(\ell\)。负指数时同一计算按极点的负阶数理解。因此必要条件是
\[
 m\mid kn,\quad m\mid\ell.
\]
这两个条件也充分：在它们成立时，全部表达式
\[
 h_\omega(\zeta)=\omega\,
       \zeta^{kn/m}(\zeta^n-1)^{\ell/m},\quad \omega^m=1
\]
都在 \(U_n\) 单值全纯且不为零，\(m\) 次方正是所需函数。任意两个提升之比取值于有限集 \(\{\omega:\omega^m=1\}\)，由 \(U_n\) 连通性必为常数，故已列出全部提升。

若 \(m\nmid\ell\)，任何 \(n\) 都不行；若 \(m\mid\ell\)，最小正整数是
\[
 n_{\min}=\frac{m}{\gcd(m,k)},
\]
其中 \(\gcd(m,0)=m\)，负 \(k\) 按其绝对值求最大公因数。第一组参数中，绕每个 \(\xi^n=1\) 的小圈始终有像绕数3，不可能被6整除，所以不存在可用的 \(n\)。第二组参数的最小值为3，全部提升为
\[
 h_\omega(\zeta)=\omega\zeta^2(\zeta^3-1),\quad \omega^6=1.
\]
幂换元把绕0的次数乘以 \(n\)，却在1的各个原像附近仍是局部一一映射，不能改变这些小圈所检测的整数 \(\ell\)。多个孔必须分别检查。
\end{solution}

\Needspace{6\baselineskip}
\begin{exercise}\label{b3:ex:11-3}
设 \(h>0\)、\(\alpha>0\)，令 \(S_h=\{w:0<\operatorname{Re}w<h\}\)。考虑由两平移 \(w\mapsto w+2\pi\mathrm i\) 与 \(w\mapsto w+2\pi\mathrm i\alpha\) 生成的群
\[
 \Gamma_\alpha=\{\,w\mapsto w+2\pi\mathrm i(m+n\alpha)
                      :m,n\in\mathbb Z\,\}.
\]
证明当 \(\alpha=p/q\) 为最简正有理数时，商是圆环，并求商坐标及模量；当 \(\alpha\) 无理时，证明作用仍自由，但不适当不连续，而且商空间不是 Hausdorff。须说明轨道为何不闭，不能只凭“群有无穷多个元素”作出判断。
\end{exercise}
\begin{solution}
若 \(\alpha=p/q\)、\(p,q\) 互素，则整数线性组合 \(mq+np\) 遍历全部整数，所以平移量组成 \((2\pi\mathrm i/q)\mathbb Z\)。映射
\[
 z=e^{qw}
\]
把 \(S_h\) 映到 \(1<|z|<e^{qh}\)，两点的像相同当且仅当相差所给平移量。局部对数给逆，故它诱导商与该圆环的双全纯对应，模量为 \(qh/(2\pi)\)。平移量有正的最小非零模长，也保证每点有与非恒等平移像不交的小邻域，紧集之间只有有限次相交。

设 \(\alpha\) 无理。将 \(0,\alpha,\ldots,M\alpha\) 的小数部分分入 \(M\) 个等长区间，抽屉原理给非零整数线性组合 \(m+n\alpha\)，其绝对值小于 \(1/M\)。必要时变号，便得到任意小的正数 \(\delta\in\mathbb Z+\alpha\mathbb Z\)。取一列这样的 \(\delta_j\to0\)，对应非恒等平移趋于恒等。任一点的小邻域都会与充分小的平移像相交，故作用不适当不连续；各个非恒等平移本身却没有固定点，所以作用仍自由。

进一步，\(\mathbb Z+\alpha\mathbb Z\) 在实轴稠密。给定 \(t\in\mathbb R\)，取整数 \(k_j=\lfloor t/\delta_j\rfloor\)，便有 \(|t-k_j\delta_j|<\delta_j\)，而 \(k_j\delta_j\) 仍在该加法群中。对任意 \(w\in S_h\)，其轨道于是稠密于经过 \(w\) 的整条竖直线。加法群可数，可取 \(t\) 不属于它；点 \(w+2\pi\mathrm i t\) 不在原轨道，却在轨道闭包中。若商是 Hausdorff，商中单点为闭集，它在商投影下的原像也应闭；这个原像正是上述不闭的轨道，矛盾。自由性本身不足以产生曲面商。
\end{solution}

\Needspace{6\baselineskip}
\begin{exercise}\label{b3:ex:11-4}
设 \(X,Y\) 为连通、Hausdorff、第二可数的曲面，\(p:X\to Y\) 是覆盖映射。证明 \(p\) 适当，即每个紧集的原像紧，当且仅当覆盖层数有限。由此证明在 \(Y\) 紧时，\(X\) 紧当且仅当覆盖有限层；并给出无穷层全纯覆盖且底曲面紧、源曲面不紧的实例。证明中须说明从何处取得闭包紧含均匀覆盖邻域的小邻域。
\end{exercise}
\begin{solution}
先设覆盖为 \(n\) 层、\(n<\infty\)。给 \(y\in Y\)，取均匀覆盖邻域 \(V\)。曲面坐标可再选在 \(V\) 内；在该坐标图中取一个较小欧氏圆盘，其闭圆盘仍含于图内。将这个小圆盘拉回，得到开邻域 \(W\)，其闭包 \(\overline W\) 是 \(V\) 内的紧集。这里使用的是曲面的局部欧氏结构；Hausdorff 性保证所取紧闭圆盘在曲面中仍为闭集。

\(V\) 上恰有 \(n\) 层，每一层的逆同胚把 \(\overline W\) 映成一个紧集，所以 \(p^{-1}(\overline W)\) 是有限个紧集的并。给任意紧集 \(K\subset Y\)，从这些 \(W\) 中选有限个覆盖 \(K\)。则 \(p^{-1}(K)\) 包含于相应有限个 \(p^{-1}(\overline W)\) 的并，且因 \(K\) 在 Hausdorff 空间 \(Y\) 中闭而为闭集。因此 \(p^{-1}(K)\) 紧，证明 \(p\) 适当。

反向若 \(p\) 适当，任意纤维 \(p^{-1}(y)\) 是紧集，又由覆盖定义是离散空间。以各个单点开集覆盖它，紧性迫使纤维有限。连通底曲面上的层数固定，故覆盖为有限层。若 \(Y\) 紧，有限层情形由适当性直接得到 \(X=p^{-1}(Y)\) 紧；若 \(X\) 紧，纤维是 \(X\) 中闭的离散子空间，同样必须有限。

复曲面的反例是
\[
 \mathbb C\longrightarrow\mathbb C/(\mathbb Z+\mathrm i\mathbb Z).
\]
格点平移自由；对两个紧集，能使它们相交的平移量落在一个有界集合中，其中格点只有有限个，故作用适当不连续。正文的商构造使投影成为全纯覆盖。闭单位方形 \(\{x+\mathrm i y:0\le x,y\le1\}\) 的像覆盖整个商，所以底曲面紧；纤维是格点集的平移，具有无穷多个点，而源 \(\mathbb C\) 不紧。这也表明紧目标本身不能代替有限层条件。
\end{solution}

\Needspace{6\baselineskip}
\begin{exercise}\label{b3:ex:11-5}
设 \(p:X\to Y\) 是连通 Riemann 曲面之间的 \(n\) 层全纯覆盖，\(n<\infty\)。证明 \(|\operatorname{Deck}(p)|\) 整除 \(n\)，并且等于 \(n\) 当且仅当覆盖正规。再令
\[
 P(w)=w^6-3w^2,\quad B=\{-2,0,2\},\quad
 X=\mathbb C\setminus P^{-1}(B),\quad Y=\mathbb C\setminus B.
\]
证明 \(P:X\to Y\) 是六层覆盖，其甲板群恰为 \(\{w\mapsto w,\ w\mapsto-w\}\)，从而给出甲板群阶数既不为1也不等于层数的例子。确定全部甲板变换时应从 \(P\circ T=P\) 出发，不预设源自同构必为线性函数。
\end{exercise}
\begin{solution}
固定纤维中的一点 \(x_0\)。若两个甲板变换在 \(x_0\) 取相同值，它们之商便固定 \(x_0\)，由 \cref{b3:prop:11-deck-free}只能为恒等。因此取值映射 \(T\mapsto T(x_0)\) 是到 \(n\) 点纤维的单射，甲板群有限。它在纤维上自由作用，每个轨道有 \(|\operatorname{Deck}(p)|\) 个点；轨道分拆整个纤维，所以群阶整除 \(n\)。等号恰好意味着只有一个轨道，也就是正规覆盖的传递性条件。

对给定多项式，
\[
 P'(w)=6w(w^4-1).
\]
临界点0给临界值0；\(w=\pm1\) 给 \(-2\)，\(w=\pm\mathrm i\) 给2，故临界值恰为 \(B\)。两个去点平面都是连通 Riemann 曲面。多项式的最高次项保证它在平面上适当，删去临界值及其全部原像后仍适当且无临界点；\cref{b3:prop:11-proper-unbranched}因此给覆盖。对每个 \(y\notin B\)，方程 \(P(w)=y\) 有六个互异根，故为六层。

设 \(T\) 是该覆盖的任意甲板变换。它在 \(X\) 全纯，并满足 \(P(T(w))=P(w)\)。最高次增长给常数 \(C_0>0\)，使
\[
 |u|\le C_0\bigl(1+|P(u)|\bigr)^{1/6}
 \quad(u\in\mathbb C).
\]
在每个被删去的有限点附近，\(P(w)\) 有界，故上述恒等式使 \(T(w)\) 有界。可去奇点定理将 \(T\) 延为整函数，恒等式也继续成立。再由 \(|P(w)|\le |w|^6+3|w|^2\)，可增大常数而得
\[
 |T(w)|\le C(1+|w|),\quad w\in\mathbb C.
\]
在任意中心 \(w_0\) 的半径 \(R\) 圆周上使用 Cauchy 估计，有
\[
 |T''(w_0)|\le\frac{2C(1+|w_0|+R)}{R^2}
 \longrightarrow0.
\]
因此 \(T(w)=aw+b\)，其中 \(a\ne0\)。

将它代入 \((aw+b)^6-3(aw+b)^2=w^6-3w^2\)。最高项给 \(a^6=1\)，\(w^5\) 项给 \(6a^5b=0\)，故 \(b=0\)；\(w^2\) 项再给 \(a^2=1\)。所以只能有 \(T(w)=w\) 或 \(T(w)=-w\)。两者确实保持 \(P\) 与源域，且都是自同构，故甲板群恰有两个元素。六点纤维分为三个甲板轨道，覆盖不正规。纤维有六层，并不意味着能够以全局对称任意交换六层。
\end{solution}

\Needspace{6\baselineskip}
\begin{exercise}\label{b3:ex:11-6}
设 \(m\ge2\)，\(p_m(w)=w^m:\mathbb C^*\to\mathbb C^*\)，且 \(F\) 在 \(\mathbb C^*\) 全纯，Laurent 展开为 \(F(w)=\sum_{n\in\mathbb Z}a_nw^n\)。分别用系数条件刻画函数 \(F\) 和微分 \(F(w)\,\mathrm dw\) 能否下降，并写出下降后的表达式。这里微分下降的含义是存在 \(\mathbb C^*\) 上全纯 \(g\)，使
\[
 F(w)=m w^{m-1}g(w^m),
\]
从而 \(F(w)\,\mathrm dw\) 是 \(g(z)\,\mathrm dz\) 的换元表达。若微分能下降，再判断其在底空间何时有单值全纯原函数。对 \(m=3\)，分别检验 \(F_1(w)=w^5+w^{-1}\) 与 \(F_2(w)=w^6+w^{-3}\)。
\end{exercise}
\begin{solution}
令 \(\zeta=e^{2\pi\mathrm i/m}\)。函数下降当且仅当 \(F(\zeta w)=F(w)\)，因为幂覆盖的纤维恰为这个旋转群的轨道。Laurent 系数唯一性把条件改写为
\[
 (\zeta^n-1)a_n=0\quad(n\in\mathbb Z).
\]
所以只有 \(m\mid n\) 的系数可以非零；下降函数为
\[
 f(z)=\sum_{j\in\mathbb Z}a_{mj}z^j.
\]
这也可先由纤维常值性作全纯下降，再比较 Laurent 展开得到。

微分则还包含变量微分的变化。题设等式若成立，便有 \(\zeta F(\zeta w)=F(w)\)。反向在这个条件下，
\[
 \frac{F(w)}{m w^{m-1}}
\]
在每个纤维常值，所以能下降为全纯 \(g\)。系数条件因此是 \((\zeta^{n+1}-1)a_n=0\)，即只有 \(m\mid n+1\) 的项可出现。写 \(n=mj-1\)，所求表达式为
\[
 g(z)=\frac1m\sum_{j\in\mathbb Z}a_{mj-1}z^{j-1}.
\]
这些级数在相应紧圆环上一致收敛；也可将它们看作已经由局部逆构造的下降函数的 Laurent 展开。\(g\) 的留数系数为 \(a_{-1}/m\)，故由圆环周期判别，底微分有单值全纯原函数当且仅当 \(a_{-1}=0\)。

在 \(m=3\) 时，\(F_1\) 的两个指数5、\(-1\) 都不被3整除，故函数不下降；但它们各加1后都被3整除，所以微分下降为
\[
 \frac13(z+z^{-1})\,\mathrm dz.
\]
其留数为 \(1/3\)，在底空间没有单值原函数。\(F_2\) 的指数6、\(-3\) 都被3整除，函数下降为 \(z^2+z^{-1}\)；这些指数加1后均不被3整除，故 \(F_2(w)\,\mathrm dw\) 不能下降。函数不变与微分不变使用不同的变换规则，不能省去微分带来的因子。
\end{solution}

\Needspace{6\baselineskip}
\begin{exercise}\label{b3:ex:11-7}
在上半平面令
\[
 T(z)=z+1,\quad U(z)=\frac{z}{z+1},\quad S(z)=-\frac1z.
\]
把 \(U\) 写成 \(S,T\) 的乘积，判断 \(T,U,TU,TU^{-1}\) 各自的类型，并求 \(TU^{-1}\) 的阶及上半平面不动点。这里乘积按映射复合解释。说明：即使两个生成元各自都没有内部不动点，它们生成的群仍未必能作为普通覆盖 \(\mathbb H\to Y\) 的甲板变换群。
\end{exercise}
\begin{solution}
使用行列式1的代表矩阵，
\[
 T=\begin{pmatrix}1&1\\0&1\end{pmatrix},\quad
 U=\begin{pmatrix}1&0\\1&1\end{pmatrix},\quad
 S=\begin{pmatrix}0&-1\\1&0\end{pmatrix}.
\]
直接相乘得 \(U=ST^{-1}S^{-1}\)。\(T\) 与 \(U\) 都是非恒等且迹为2的抛物元，边界不动点分别是 \(\infty\) 与0。它们的全部非恒等幂也没有内部固定点。

乘积的矩阵为
\[
 TU=\begin{pmatrix}2&1\\1&1\end{pmatrix},\quad
 R\coloneqq TU^{-1}=\begin{pmatrix}0&1\\-1&1\end{pmatrix}.
\]
第一矩阵的迹为3，故为双曲型，两个边界不动点是 \(z^2-z-1=0\) 的根 \((1\pm\sqrt5)/2\)。第二矩阵的迹为1，故为椭圆型；它的分式是 \(R(z)=1/(1-z)\)，不动点方程为 \(z^2-z+1=0\)，上半平面根为
\[
 z_0=\frac{1+\sqrt3\,\mathrm i}{2}.
\]
还可直接计算
\[
 R^2=\begin{pmatrix}-1&1\\-1&0\end{pmatrix},\quad R^3=-I.
\]
在 \(\operatorname{PSL}_2(\mathbb R)\) 中 \(-I\) 与恒等相同，且 \(R,R^2\) 均非恒等，所以 \(R\) 恰有三阶。

因此由 \(T,U\) 生成的群含一个非恒等而固定 \(z_0\) 的变换。连通覆盖的甲板变换一旦固定一点就必须恒等，所以这组给定变换不能成为普通覆盖 \(\mathbb H\to Y\) 的甲板群。两个生成元分别是抛物元，只检查生成元并没有排除其乘积产生椭圆元；自由性是对整个群的要求。
\end{solution}

\Needspace{6\baselineskip}
\begin{exercise}\label{b3:ex:11-8}
设 \(a>1\)，在 \(\mathbb H\) 上令 \(D(z)=az\)、\(R(z)=-1/z\)。可使用正文的缩放商模型
\[
 q(z)=\exp\left(\frac{2\pi\mathrm i\log z}{\log a}\right),\quad
 r=\exp\left(-\frac{2\pi^2}{\log a}\right),\quad
 \mathbb H/\langle D\rangle\simeq A(r,1),
\]
其中 \(0<\operatorname{Im}\log z<\pi\)。证明 \(R\) 诱导圆环上的自同构，求其公式与固定点。再证明
\[
 J(\zeta)=\zeta+\frac r\zeta
\]
把这个圆环满射到椭圆内部
\[
 E=\left\{\,x+\mathrm i y:
       \frac{x^2}{(1+r)^2}+\frac{y^2}{(1-r)^2}<1\,\right\},
\]
其纤维恰为上述自同构的轨道。求分歧点与分歧值，说明为何在缩放群中加入 \(R\) 后，商投影不能继续称为普通无分歧覆盖。
\end{exercise}
\begin{solution}
因为 \(RDR^{-1}=D^{-1}\)，\(R\) 将每个缩放轨道送到另一个缩放轨道，因而可在商上定义映射。取所给对数分支，有
\[
 \log(-1/z)=-\log z+\pi\mathrm i.
\]
代入 \(q\) 即得 \(q(Rz)=r/q(z)\)，所以诱导的自同构是
\[
 \iota(\zeta)=\frac r\zeta.
\]
它将 \(r<|\zeta|<1\) 保持不变，且平方为恒等。固定点方程为 \(\zeta^2=r\)，故恰有 \(\sqrt r,-\sqrt r\) 两点，均在圆环内部。

\(J\) 在 \(\iota\) 下不变。写 \(\zeta=\rho e^{\mathrm i t}\)，则
\[
 J(\zeta)=\left(\rho+\frac r\rho\right)\cos t
       +\mathrm i\left(\rho-\frac r\rho\right)\sin t.
\]
若 \(r<\rho<1\)，由 \((\rho-r)(\rho-1)<0\) 得 \(\rho+r/\rho<1+r\)；函数 \(\rho-r/\rho\) 严格递增，其端点值为 \(-(1-r)\) 与 \(1-r\)，故其绝对值小于 \(1-r\)。代入椭圆的不等式，得到 \(J(A(r,1))\subset E\)。在两条边界 \(\rho=1\) 和 \(\rho=r\) 上，上式分别以相反方向描出同一条椭圆边界。

这还给出 \(J:A(r,1)\to E\) 的适当性。任取紧集 \(K\subset E\)，在闭圆环中的原像是紧集；它不能碰到两条边界，因为边界的像都在 \(\partial E\)。所以全部原像紧含于开圆环。非常数全纯映射的像是开集；适当性又使像在 \(E\) 中闭：若 \(J(\zeta_n)\to u\in E\)，这些像点连同极限组成紧集，原像紧性给一个收敛子列，其极限仍在圆环且映到 \(u\)。\(E\) 连通，故非空的像就是整个 \(E\)。

若 \(J(\eta)=J(\zeta)\)，两边乘以 \(\eta\zeta\) 并整理，得到
\[
 (\eta-\zeta)(\eta\zeta-r)=0.
\]
因此 \(\eta=\zeta\) 或 \(\eta=r/\zeta\)，纤维恰为 \(\iota\) 的轨道。由于 \(J\) 是开映射，它诱导轨道商与 \(E\) 的同胚，局部表达则给相应的复坐标。

导数 \(J'(\zeta)=1-r/\zeta^2\) 只在 \(\zeta_0=\pm\sqrt r\) 为零，分歧值为 \(\pm2\sqrt r\)。在这两点附近，恒等式
\[
 J(\zeta)-2\zeta_0=\frac{(\zeta-\zeta_0)^2}{\zeta}
\]
表明局部次数为2：非零因子 \(1/\zeta\) 在小盘内有全纯平方根，可吸收到局部坐标中。其余点的导数非零，普通纤维有两个点；删去两个分歧值及其原像后得到二层全纯覆盖。完整映射则在两个特殊纤维只剩一个点，并有局部平方形式。缩放作用原本自由，加入 \(R\) 后出现固定点，最终的商虽仍有明确的曲面模型，投影却已经发生分歧。
\end{solution}
